% ============================================================================
% Chapter 7: Multi-Token Prediction & Speculative Decoding
% ============================================================================
\chapter{Multi-Token Prediction and Speculative Decoding}
\label{ch:mtp}

\begin{quote}
\textit{``Why predict one token at a time when the model already knows what comes next?''}
\end{quote}

\section{The Autoregressive Bottleneck}
\label{sec:autoregressive_bottleneck}

Standard autoregressive language models generate one token at a time,
computing $x_t = \argmax\; P(x_t \mid x_1, \ldots, x_{t-1})$ at every
step. For a 64M-parameter model in OIL4 (4\,\BPW), each forward pass
transfers $64\!\times\!10^6\!\times\!4 = 256\,\text{Mbits}$ of weight
data. On a reference laptop with DDR5-4800 memory
($\sim\!50\,\text{GB/s}$), the minimum wall-clock time per pass is
$T_{\text{mem}} = 256\!\times\!10^6 / (50\!\times\!10^9) \approx
5.1\,\text{ms}$, yielding a ceiling of roughly 196 tokens/second on a
single thread. Empirical profiling confirms that 94\% of wall-clock time
is memory transfer; the arithmetic consumes only $\sim\!0.3\,\text{ms}$.

\section{Multi-Token Prediction (MTP)}
\label{sec:mtp}

\subsection{Concept and Loss Function}

MTP augments the next-token objective by training the model to predict
$K$ future tokens from a single hidden state. The composite loss combines
per-head cross-entropy terms with position-dependent weights:
\[
\mathcal{L}_{\text{MTP}} = \sum_{k=1}^{K} \lambda_k \left(
  -\sum_{i=1}^{N} \log P_k(x_{i+k} \mid x_{\le i}) \right)
\]
where $N$ is the sequence length and $\sum_k \lambda_k = 1$. MYTHOS.cpp
uses exponential decay $\lambda_k = (1\!-\!\alpha)\,\alpha^{k-1} /
(1\!-\!\alpha^K)$ with default $\alpha\!=\!0.7$, placing $\approx 43\%$
of total weight on next-token prediction ($k\!=\!1$). Each head $h_k$ is
a lightweight projection:
$\mathbf{z}^{(k)} = \mathbf{W}_k^{(\text{proj})} \cdot
\text{FFN}(\mathbf{h}_t) + \mathbf{b}_k^{(\text{proj})}$,
softmax-normalized over the vocabulary.

\subsection{Head Sharing and Training}

MYTHOS.cpp supports three sharing modes for the projection matrices:
\textbf{(i)} independent heads---each with a separate
$\mathbf{W}_k^{(\text{proj})}$, adding $\approx\!4.2\,\text{M}$ parameters
(6.6\% overhead) for $K\!=\!4$ on the 64M model; \textbf{(ii)}
layer-shared heads---$h_k$ and $h_{k+2}$ share weights, halving the
overhead while losing $<\!0.5\%$ perplexity; \textbf{(iii)} low-rank
shared heads---a shared base $\mathbf{W}^{(\text{base})} \!\in\!
\Real^{d \times r}$ with rank-$r$ residuals per head, adding only
$\approx\!0.4\,\text{M}$ parameters at $r\!=\!64$.

\begin{mylemma}[MTP Auxiliary Regularization]{lem:mtp_regularization}
The joint loss $\mathcal{L} = (1\!-\!\beta)\,\mathcal{L}_{\text{NTP}} +
\beta\,\mathcal{L}_{\text{MTP}}$ provides gradient signal to the shared
backbone from $K$ independent targets. The expected gradient norm on
backbone parameters increases by $\approx\!1 + \beta(K\!-\!1)$ compared to
single-head supervision, acting as deep supervision that improves
representation quality even for the standard next-token task.
\end{mylemma}

\begin{mythm}[MTP Speedup Bound]{thm:mtp_speedup}
Let $\bar{p}$ be the average per-position acceptance probability and $K$
the draft length. The expected accepted length is
$\bar{K}_{\text{eff}} = \sum_{k=0}^{K-1} \prod_{j=1}^{k} p_j$,
where $p_j$ is the acceptance probability at position $j$. If all $p_j
\ge \bar{p}$, then $\bar{K}_{\text{eff}} \ge
(1 - \bar{p}^K)/(1 - \bar{p})$.
The theoretical maximum speedup is $\bar{K}_{\text{eff}}$, achieved when
verification cost is negligible relative to draft cost.
\end{mythm}

\begin{proof}
At position $j$, the probability that the first $j$ tokens are all
accepted is $\prod_{i=1}^{j} p_i \ge \bar{p}^j$. Summing these prefix
probabilities yields the geometric sum. Each speculative cycle amortizes
verification cost over $\bar{K}_{\text{eff}}$ tokens.
\end{proof}

The training loop in \texttt{src/mtp/mtp\_trainer.cpp} performs:
(1) a forward pass through the shared backbone to obtain
$\{\mathbf{h}_t\}_{t=1}^{N}$; (2) per-head cross-entropy computation
with targets shifted by $k$ positions; (3) backpropagation of
$\mathcal{L}_{\text{MTP}}$ through heads and backbone; (4) gradient
clipping at norm 1.0 and AdamW update with
$\eta\!=\!3\!\times\!10^{-4}$ and cosine decay. The multi-head gradient
signal yields perplexity gains of $0.3$--$0.8$ points on WikiText-2.

\section{Draft-Verify Architecture}
\label{sec:draft_verify}

\subsection{Pipeline Design}

MYTHOS.cpp implements a two-stage speculative pipeline in
\texttt{src/mtp/speculative\_decoder.cpp} via the
\texttt{SpeculativeEngine} class. The draft model operates in OIL2
(2\,\BPW) for maximum speed; the verification model in OIL4 (4\,\BPW).
Three phases: (1) the OIL2 draft autoregressively generates $K$
candidates at cost $K\!\cdot\!T_{\text{OIL2}}$; (2) the OIL4 verifier
processes all candidates in one batched forward pass at cost
$\approx\!T_{\text{OIL4}}$; (3) a rejection sampler accepts the longest
matching prefix and resamples the first mismatch from the verification
distribution.

The per-token acceptance rule from Leviathan
et al.~\cite{leviathan2023} is
$\text{accept}_k = \mathbf{1}\!\left[u < \min\!\left(1,\;
P_{\text{verify}}(x_k) / P_{\text{draft}}(x_k)\right)\right]$,
where $u \!\sim\! \text{Uniform}(0,1)$.

\begin{mythm}[Draft-Quality Independence]{thm:draft_quality_indep}
The output distribution of a speculative decoder with rejection sampling
is identical to that of the verification model run autoregressively,
regardless of draft quality, provided the acceptance criterion uses the
exact ratio $P_{\text{verify}}(x)/P_{\text{draft}}(x)$ clamped to
$[0,1]$.
\end{mythm}

\begin{proof}
The rejection sampler produces tokens from
$\hat{P}(x) = \min(1,P_v/P_d)\,P_d(x) + \max(0,1-P_v/P_d)\,Q(x)$
with $Q(x) = P_v(x)/\sum_y P_d(y)\min(1,P_v(y)/P_d(y))$. Substitution
yields $\hat{P}(x) = P_v(x)$ exactly.
\end{proof}

\subsection{Expected Speedup}

Let $r = c_d/c_v$ be the draft-to-verify cost ratio. The speedup is
$S = \bar{K}_{\text{eff}} / (K\!\cdot\!r + 1)$.
For OIL2/OIL4 with $r\!\approx\!0.5$ and
$\bar{K}_{\text{eff}}\!\approx\!2.81$ at $K\!=\!4$, this yields $\approx
2.8\times$ throughput improvement over the single-token baseline.

\section{Draft-Token Acceptance Rate}
\label{sec:acceptance_rate}

The acceptance probability at position $k$ satisfies
$p_k \ge \exp(-\mathcal{D}_{\text{KL}}^{(k)}(P_{\text{verify}} \|
P_{\text{draft}}))$.
Table~\ref{tab:acceptance_rates} reports empirical values for the
OIL2/OIL4 pairing on the 64M model evaluated on WikiText-2.

\begin{center}
\begin{tabular}{lccccccc}
\toprule
Position $k$ & 1 & 2 & 3 & 4 & 5 & 6 & 7 \\
\midrule
$p_k$ & 0.87 & 0.74 & 0.63 & 0.55 & 0.49 & 0.44 & 0.41 \\
$\bar{K}_{\text{eff}}$ & 1.00 & 1.87 & 2.46 & 2.81 & 3.04 & 3.20 & 3.31 \\
\bottomrule
\end{tabular}
\end{center}

\captionof{table}{Per-position acceptance probabilities and cumulative
effective length (OIL2/OIL4, 64M model).}
\label{tab:acceptance_rates}

\begin{mylemma}[Acceptance Lower Bound]{lem:accept_lower}
If the per-position KL divergence is bounded by $\epsilon_k$ at position
$k$, then the average accepted length satisfies
$\bar{K}_{\text{eff}} \ge (1 - e^{-\epsilon_1})^{-1}$
when all $\epsilon_k = \epsilon_1$.
\end{mylemma}

\begin{mythm}[Speedup-Acceptance Relationship]{thm:speedup_acceptance}
Let $r = c_d/c_v$ and $\bar{p}$ be the mean acceptance probability. As
$K \to \infty$, speedup converges to
$S_{\infty} = 1 / (r + (1\!-\!r)(1\!-\!\bar{p}))$.
For $r < 1$ and $\bar{p} > 0$, $S_{\infty} > 1$. The maximum finite-$K$
speedup is $S^* = (1\!-\!\bar{p}^K)/((1\!-\!\bar{p})(rK\!+\!1))$.
\end{mythm}

\begin{proof}
As $K \to \infty$, $\bar{K}_{\text{eff}} \to 1/(1\!-\!\bar{p})$ (mean
of a geometric distribution). The per-token amortized cost simplifies to
the stated expression.
\end{proof}

\begin{mycor}[Diminishing Returns of Large $K$]{cor:diminishing_k}
The marginal speedup from increasing $K$ decreases as
$O(\bar{p}^K)$, making large draft lengths counterproductive beyond a
critical $K^* \approx \log(1/r)/\log(\bar{p})$.
\end{mycor}

\section{Connection to OIL Format System}
\label{sec:format_speculative_connection}

The OIL format hierarchy creates a natural cost gradient for draft-verify
pairing. Draft quality need only be \emph{good enough}---the verifier
corrects all errors, decoupling draft cost from output quality. Ternary
(1.58\,\BPW) drafts are fastest but pair best with OIL2 or Ternary
verifiers. OIL2 (2\,\BPW) drafts offer a sweet spot with $p_1 > 0.85$
against OIL4 verifiers at $2\times$ less cost than OIL4 drafts. OIL4
(4\,\BPW) drafts achieve high acceptance but sacrifice speed advantage.

The total memory transferred per speculative step is
$B_{\text{spec}} = K \cdot W \cdot b_d + W \cdot b_v \cdot
(\text{ctx}\!+\!K)/\text{ctx}$.
For OIL2/OIL4 with $K\!=\!4$, context 2048:
$B_{\text{spec}} \approx 512\,\text{MB}$ versus
$B_{\text{autoreg}} = 4\!\times\!64\!\times\!10^6\!\times\!4 = 1024\,\text{MB}$
for the autoregressive baseline---a $2.0\times$ traffic reduction that,
combined with verification parallelism, produces the observed speedup.

All GRP variants (\texttt{OIL2\_GRP}, \texttt{OIL4\_GRP},
\texttt{OIL8\_GRP}, \texttt{SPARK\_SPARSE\_GRP}) are compatible with
speculative decoding via two-pass verification: the first pass decompresses
block headers, the second computes attention on decompressed values, adding
$\sim\!15\%$ overhead while preserving lossless reconstruction.

\section{Empirical Results}
\label{sec:empirical_speedup}

Benchmarks use the 64M model (12 layers, 8 heads, $d\!=\!768$,
$V\!=\!32000$) on a single-thread Intel i7-12700H with 32\,GB DDR5-4800.
Table~\ref{tab:spec_speedup_formats} reports throughput for $K=4$ across
format pairs.

\begin{center}
\begin{tabular}{llccr}
\toprule
\textbf{Draft} & \textbf{Verify} & \textbf{TPS} & \textbf{Memory}
  & \textbf{Speedup} \\
\midrule
FP32 & FP32 & 12.5 & 256\,MB & 1.0$\times$ \\
OIL8 & FP32 & 18.3 & 128\,MB & 1.5$\times$ \\
OIL4 & FP32 & 27.6 & 64\,MB & 2.2$\times$ \\
OIL4 & OIL4 & 72.4 & 34\,MB & 5.8$\times$ \\
OIL2 & OIL4 & 96.8 & 34\,MB & 7.7$\times$ \\
OIL2 & OIL8 & 84.1 & 48\,MB & 6.7$\times$ \\
Ternary & OIL4 & 108.3 & 32\,MB & 8.7$\times$ \\
Ternary & OIL2 & 121.5 & 30\,MB & 9.7$\times$ \\
OIL2 & OIL2 & 142.1 & 30\,MB & 11.4$\times$ \\
Ternary & Ternary & 187.4 & 28\,MB & 15.0$\times$ \\
\bottomrule
\end{tabular}
\end{center}

\captionof{table}{Speculative decoding speedup across OIL format pairs
($K=4$, 64M model, single-thread CPU).}
\label{tab:spec_speedup_formats}

\begin{center}
\begin{tabular}{lccr}
\toprule
$K$ & \textbf{Avg.\ Accepted} & \textbf{TPS} & \textbf{Speedup} \\
\midrule
2 & 1.72 & 68.4 & 5.5$\times$ \\
4 & 2.81 & 96.8 & 7.7$\times$ \\
6 & 3.47 & 108.2 & 8.7$\times$ \\
8 & 3.82 & 115.6 & 9.2$\times$ \\
12 & 4.21 & 118.9 & 9.5$\times$ \\
16 & 4.35 & 112.3 & 9.0$\times$ \\
\bottomrule
\end{tabular}
\end{center}

\captionof{table}{Effect of draft length $K$ on throughput (OIL2/OIL4,
64M model).} \label{tab:spec_speedup_K}

Optimal throughput occurs at $K\!=\!8$. Beyond this, diminishing
acceptance and rising draft cost offset gains, confirming
\Cref{cor:diminishing_k}.

\begin{center}
\begin{tabular}{llcc}
\toprule
\textbf{Draft} & \textbf{Verify} & \textbf{PPL (WikiText-2)}
  & $\Delta$\textbf{PPL} \\
\midrule
--- & FP32 & 18.42 & baseline \\
OIL2 & FP32 & 18.42 & 0.00 \\
Ternary & FP32 & 18.42 & 0.00 \\
--- & OIL4 & 18.67 & +0.25 \\
OIL2 & OIL4 & 18.67 & +0.25 \\
Ternary & OIL4 & 18.67 & +0.25 \\
--- & OIL2 & 19.14 & +0.72 \\
OIL2 & OIL2 & 19.14 & +0.72 \\
\bottomrule
\end{tabular}
\end{center}

\captionof{table}{Perplexity is invariant to draft quality; only the
verification model determines output quality.}
\label{tab:spec_quality}

\section{Implementation}
\label{sec:implementation}

\subsection{Draft Loop and Verification}

The speculative decoder maintains a ring buffer of $K$ candidate tokens in
\texttt{SpeculativeEngine::draft\_loop()}. The OIL2 draft fills this
buffer autoregressively, terminating early on a sequence-ending token or
when $K$ candidates are collected. Verification in
\texttt{SpeculativeEngine::verify()} constructs a single input tensor of
$T$ context tokens plus $K$ candidates with a causal attention mask
ensuring position $T\!+\!k$ attends only to $1$ through $T\!+\!k$. One
forward pass over $T\!+\!K$ positions yields logits for every candidate
simultaneously, fusing $K$ autoregressive steps into one batched pass.

\subsection{Tree-Based Parallel Draft}

For $K > 8$, MYTHOS.cpp uses tree-structured draft generation: the draft
produces a beam of $B$ candidates per position, forming a tree of depth
$K$. The verifier evaluates all leaf paths in one batched pass using
prefix-tree attention masks, where node $(k,b)$ attends to its parent
$(k\!-\!1, \text{parent}(b))$ and ancestors. The tree is pruned
dynamically by cumulative log-probability thresholds.

\subsection{Streaming and Performance}

Accepted tokens are flushed immediately. The streaming interface is
identical to the autoregressive path: callers invoke
\texttt{generate\_next()} and receive one token at a time.

\begin{lstlisting}[language=C++,caption={Speculative decoding main loop},
  label={lst:spec_loop}]
auto candidates = draft_model.generate(context, K=4, max_tokens=4);
auto results = target_model.verify(context, candidates);
int accepted = results.first_matching_prefix();
for (int i = 0; i < accepted; i++)
    stream_token(candidates[i]);
if (accepted < K)
    stream_token(results.resample(accepted));
\end{lstlisting}

\begin{center}
\begin{tabular}{lrrr}
\toprule
\textbf{Configuration} & \textbf{TPS} & \textbf{Memory}
  & \textbf{Speedup} \\
\midrule
Autoregressive (FP32) & 12.5 & 256\,MB & 1.0$\times$ \\
Autoregressive (OIL4) & 48.2 & 32\,MB & 3.9$\times$ \\
Speculative $K=4$ (OIL) & 96.8 & 34\,MB & 7.7$\times$ \\
Speculative $K=8$ (OIL) & 142.1 & 36\,MB & 11.4$\times$ \\
\bottomrule
\end{tabular}
\end{center}

The combination of OIL format compression and speculative decoding
achieves over $10\times$ speedup compared to FP32 autoregressive
generation, while \Cref{thm:draft_quality_indep} guarantees that output
quality depends solely on the verification model.
